\documentclass[11pt,a4paper]{article}
% preamble.tex — shared setup for human-verification documents
% Compile each document with:  tectonic <File>.tex

\usepackage[T1]{fontenc}
\usepackage[utf8]{inputenc}
\usepackage{lmodern}
\usepackage{amssymb,amsmath,amsthm}
\usepackage{mathtools}
\usepackage[dvipsnames]{xcolor}
\usepackage{listings}
\usepackage{hyperref}
\usepackage[margin=2.5cm]{geometry}
\usepackage{parskip}

% ── Colours ──────────────────────────────────────────────────────────
\definecolor{leanBg}{HTML}{F7F7F0}
\definecolor{haskellBg}{HTML}{F0F4F7}
\definecolor{keyword}{HTML}{7B3294}
\definecolor{comment}{HTML}{5A8C32}
\definecolor{noteFg}{HTML}{B35900}

% ── Lean 4 listings language ─────────────────────────────────────────
% NOTE: Lean listing blocks use ASCII approximations of unicode symbols
% (e.g. -> for arrow, o. for compose).  The literate key substitutes
% these with proper math-mode renderings in the PDF.
\lstdefinelanguage{Lean4}{
  morekeywords={def,theorem,structure,where,namespace,end,import,fun,
                let,have,show,by,exact,intro,cases,rfl,subst,obtain,
                funext,simp,rw,apply,variable,Type,Prop,instance,
                section,open,class,extends,deriving,sorry,noncomputable},
  sensitive=true,
  morecomment=[l]{--},
  morecomment=[n]{/-}{-/},
  morestring=[b]",
}

% ── Lean listing style ───────────────────────────────────────────────
\lstdefinestyle{lean}{
  language=Lean4,
  backgroundcolor=\color{leanBg},
  basicstyle=\small\ttfamily,
  keywordstyle=\color{keyword}\bfseries,
  commentstyle=\color{comment}\itshape,
  breaklines=true,
  frame=single,
  framerule=0.4pt,
  xleftmargin=4pt,
  xrightmargin=4pt,
  aboveskip=6pt,
  belowskip=4pt,
  literate=
    {->}{{$\to$}}2
    {=>}{{$\Rightarrow$}}2
    {<|}{{$\langle$}}1
    {|>}{{$\rangle$}}1
    {==>}{{$\Longrightarrow$}}3
    { o. }{{$\;\circ\;$}}3
    {forall}{{$\forall$}}1
    {:=}{{$\,:=\,$}}2
    {!=}{{$\neq$}}2
    {<=}{{$\leq$}}2
    {>=}{{$\geq$}}2,
}

% ── Haskell listing style ───────────────────────────────────────────
\lstdefinestyle{haskell}{
  language=Haskell,
  backgroundcolor=\color{haskellBg},
  basicstyle=\small\ttfamily,
  keywordstyle=\color{keyword}\bfseries,
  commentstyle=\color{comment}\itshape,
  breaklines=true,
  frame=single,
  framerule=0.4pt,
  xleftmargin=4pt,
  xrightmargin=4pt,
  aboveskip=6pt,
  belowskip=4pt,
  literate=
    {->}{{$\to$}}2
    {=>}{{$\Rightarrow$}}2,
}

% ── Theorem environments ────────────────────────────────────────────
\theoremstyle{definition}
\newtheorem{defn}{Definition}[section]
\newtheorem{thm}[defn]{Theorem}
\newtheorem{prop}[defn]{Proposition}

\theoremstyle{remark}
\newtheorem*{fidelitynote}{Fidelity note}

% ── Convenience macros ───────────────────────────────────────────────
\newcommand{\lean}[1]{\texttt{#1}}      % inline Lean identifier
\newcommand{\hask}[1]{\texttt{#1}}      % inline Haskell identifier
\newcommand{\Ext}{\mathrm{Ext}}
\newcommand{\Pos}{\mathrm{Pos}}
\newcommand{\Shape}{\mathrm{Shape}}


\title{Human Verification: \texttt{Containers.Sourcing}\\[4pt]
  \large Sourcing as a Functor from Routed BOMs to the Opposite Network}
\author{Machine-generated from \texttt{lean/Sourcing.lean} for review by Neil Ghani}
\date{July 2026}

\begin{document}
\maketitle

\begin{abstract}
This document presents every definition and theorem \emph{statement} from the
Lean~4 module \lean{Containers.Sourcing} in three parallel forms: the original
Lean, a readable Haskell rendering, and precise mathematical English.  Proofs
are omitted --- they are machine-checked by Lean.  The purpose is to let a
human verifier confirm that the formalisation captures the intended
mathematics.

The file verifies Construction~\texttt{con:sourcing} (\S6 of the paper): the
sourcing assignment --- ``each product is made somewhere, each component
occurrence travels an inbound route to its point of use'' --- is a covariant
functor of directed containers from the directed container of
\emph{routed BOMs} to the \emph{opposite shipping network}.  The file ends
with the UAE steel rebar BOM as a worked example.

A formalization caveat (flagged to Robin and Neil): the paper's domain is
\emph{all} product-labelled trees, but for a general tree the sourcing
position map requires routes that may not exist.  The fix adopted here is to
carry routing choices \emph{on the shapes}: a routed BOM bundles each BOM
edge with the inbound route of its component.
\end{abstract}

\tableofcontents
\bigskip

%===================================================================
\section{Overview}
%===================================================================

The \textbf{opposite digraph} $G^{\mathrm{op}}$ of $G$ swaps source and
target: a path in $G^{\mathrm{op}}$ from $v$ to $u$ is a path in $G$ from $u$
to $v$.  Positions of $\mathcal{N}(G^{\mathrm{op}})$ at $v$ are
\emph{inbound routes} --- paths of $G$ ending at $v$.

A \textbf{routed BOM} over products $P$ and graph $G$, indexed by its
producing facility $v$, is a product-labelled tree whose every edge carries
the inbound route (in $G^{\mathrm{op}}$) of that component to its parent.
Occurrences navigate this tree; the sourcing functor maps each occurrence to
its total inbound route (by appending the routes along the path from root to
occurrence).


%===================================================================
\section{The opposite digraph}
%===================================================================

% --- Digraph.op ----------------------------------------------------
\subsection{Definition: \lean{Digraph.op}}

\begin{defn}[Opposite digraph]
The \textbf{opposite} of a digraph $G$ is the digraph $G^{\mathrm{op}}$ with
the same nodes and edges, but with source and target swapped:
$G^{\mathrm{op}}.\mathrm{src} = G.\mathrm{tgt}$ and
$G^{\mathrm{op}}.\mathrm{tgt} = G.\mathrm{src}$.

Positions of $\mathcal{N}(G^{\mathrm{op}})$ at $v$ are pairs $(u, \rho)$
where $\rho$ is a path in $G$ from $u$ to $v$ --- the \emph{inbound routes}
arriving at $v$.
\end{defn}

\paragraph{Lean.}
\begin{lstlisting}[style=lean]
@[reducible] def Digraph.op (G : Digraph) : Digraph where
  V := G.V
  E := G.E
  src := G.tgt
  tgt := G.src
\end{lstlisting}

\paragraph{Haskell.}
\begin{lstlisting}[style=haskell]
op :: Digraph v e -> Digraph v e
op g = Digraph { src = tgt g, tgt = src g }
-- A path in (op g) from v to u is a G-path from u to v.
\end{lstlisting}


%===================================================================
\section{Routed BOMs}
%===================================================================

% --- Bom -----------------------------------------------------------
\subsection{Inductive type: \lean{Bom}}

\begin{defn}[Routed BOM tree]
A \textbf{routed BOM} $\lean{Bom}(P, G, v)$ over products $P$ and graph $G$,
indexed by producing facility $v$, is an inductive tree:
\[
  \lean{node}(\mathrm{prod},\, v,\, n,\, \mathrm{sites},\, \mathrm{routes},\, \mathrm{children})
\]
where:
\begin{itemize}
  \item $\mathrm{prod} : P$ is the product label at this node;
  \item $v : G.V$ is the facility producing \lean{prod};
  \item $n : \mathbb{N}$ is the number of direct components;
  \item $\mathrm{sites} : \mathrm{Fin}(n) \to G.V$ gives the producing facility
    of each component;
  \item $\mathrm{routes}(i) : \mathrm{Path}_{G^{\mathrm{op}}}(v, \mathrm{sites}(i))$
    is the inbound route of component $i$ (a $G$-path from $\mathrm{sites}(i)$
    to $v$, reversed);
  \item $\mathrm{children}(i) : \lean{Bom}(P, G, \mathrm{sites}(i))$ is the
    sub-BOM for component $i$.
\end{itemize}
Raw materials are nodes with $n = 0$ (no children).
\end{defn}

\paragraph{Lean.}
\begin{lstlisting}[style=lean]
inductive Bom (P : Type) (G : Digraph) : G.V -> Type where
  | node (prod : P) (v : G.V) (n : Nat) (sites : Fin n -> G.V)
      (routes : (i : Fin n) -> Digraph.Path G.op v (sites i))
      (children : (i : Fin n) -> Bom P G (sites i)) : Bom P G v
\end{lstlisting}

\paragraph{Haskell.}
\begin{lstlisting}[style=haskell]
data Bom p v e = Node
  { prodLabel :: p
  , site      :: v
  , children  :: [(v, [e], Bom p v e)]
    -- each child: (site, inbound route, sub-BOM)
  }
\end{lstlisting}

\begin{fidelitynote}
In Lean, \lean{Bom P G v} is indexed by the producing facility $v$, so the
type system enforces that routes and children are consistently sited.  In
Haskell, $v$ is a type parameter (not an index), so this constraint is a
runtime property only.
\end{fidelitynote}


% --- Bom.Occ -------------------------------------------------------
\subsection{Inductive type: \lean{Bom.Occ}}

\begin{defn}[Occurrence]
An \textbf{occurrence} in a routed BOM $t : \lean{Bom}(P, G, v)$ is a
descending path from the root:
\begin{itemize}
  \item $\lean{here} : \lean{Occ}(t)$ --- the root occurrence;
  \item $\lean{inside}(i, p) : \lean{Occ}(t)$ where $t =
    \lean{node}(\ldots, \mathrm{children})$, $i : \mathrm{Fin}(n)$, and
    $p : \lean{Occ}(\mathrm{children}(i))$ --- an occurrence inside the $i$-th
    child.
\end{itemize}
\end{defn}

\paragraph{Lean.}
\begin{lstlisting}[style=lean]
inductive Occ : {v : G.V} -> Bom P G v -> Type where
  | here   {v : G.V} {t : Bom P G v} : Occ t
  | inside {prod : P} {v : G.V} {n : Nat} {sites : Fin n -> G.V}
      {routes : (i : Fin n) -> Digraph.Path G.op v (sites i)}
      {children : (i : Fin n) -> Bom P G (sites i)}
      (i : Fin n) (p : Occ (children i)) :
      Occ (.node prod v n sites routes children)
\end{lstlisting}

\paragraph{Haskell.}
\begin{lstlisting}[style=haskell]
data Occ = Here | Inside Int Occ
  -- Here: root occurrence
  -- Inside i p: occurrence p in the i-th child
\end{lstlisting}


% --- Bom.sub / Bom.shift ------------------------------------------
\subsection{Definitions: \lean{Bom.sub} and \lean{Bom.shift}}

\begin{defn}[Sub-BOM and occurrence shift]
Given $t : \lean{Bom}(P, G, v)$ and $p : \lean{Occ}(t)$:
\begin{itemize}
  \item $\lean{sub}(t, p) : (u : G.V) \times \lean{Bom}(P, G, u)$ is the
    sub-BOM at occurrence $p$ (with its producing facility).
  \item $\lean{shift}(t, p, q) : \lean{Occ}(t)$ for $q :
    \lean{Occ}(\lean{sub}(t, p).\mathbf{2})$ is the concatenation: an occurrence
    $q$ inside the sub-BOM at $p$, seen from the root of $t$.
\end{itemize}
\end{defn}

\paragraph{Lean.}
\begin{lstlisting}[style=lean]
def sub : {v : G.V} -> (t : Bom P G v) -> Occ t
    -> (u : G.V) * Bom P G u
  | v, t, .here => <|v, t|>
  | _, .node _ _ _ _ _ children, .inside i p => (children i).sub p

def shift : {v : G.V} -> (t : Bom P G v) -> (p : Occ t)
    -> Occ (t.sub p).2 -> Occ t
  | _, _, .here, q => q
  | _, .node _ _ _ _ _ children, .inside i p, q =>
      .inside i ((children i).shift p q)
\end{lstlisting}

\paragraph{Haskell.}
\begin{lstlisting}[style=haskell]
subBom :: Bom p v e -> Occ -> (v, Bom p v e)
subBom t Here = (site t, t)
subBom (Node _ _ children) (Inside i p) =
  subBom (thd (children !! i)) p

shift :: Bom p v e -> Occ -> Occ -> Occ
shift _ Here q = q
shift (Node _ _ children) (Inside i p) q =
  Inside i (shift (thd (children !! i)) p q)
\end{lstlisting}


% --- DC laws -------------------------------------------------------
\subsection{Theorems: \lean{Bom.shift\_root}, \lean{Bom.sub\_shift}, \lean{Bom.shift\_assoc}}

\begin{thm}[BOM directed-container laws]
For a routed BOM $t : \lean{Bom}(P, G, v)$:
\begin{enumerate}
  \item \textbf{D3} (\lean{shift\_root}): $\lean{shift}(t, p, \lean{here}) = p$
    for all $p$.
  \item \textbf{D4} (\lean{sub\_shift}): $\lean{sub}(t, \lean{shift}(t, p, q))
    = \lean{sub}(\lean{sub}(t, p).\mathbf{2}, q)$ (sub-BOMs nest along the shift ---
    multi-level costing is coherent).
  \item \textbf{D5} (\lean{shift\_assoc}): the shift is associative (modulo
    transport along D4).
\end{enumerate}
\end{thm}

\paragraph{Lean.}
\begin{lstlisting}[style=lean]
theorem shift_root : {v : G.V} -> (t : Bom P G v) -> (p : Occ t) ->
    t.shift p .here = p

theorem sub_shift : {v : G.V} -> (t : Bom P G v) -> (p : Occ t) ->
    (q : Occ (t.sub p).2) -> t.sub (t.shift p q) = (t.sub p).2.sub q

theorem shift_assoc : {v : G.V} -> (t : Bom P G v) -> (p : Occ t) ->
    (q : Occ (t.sub p).2) ->
    (r : Occ ((t.sub p).2.sub q).2) ->
    t.shift (t.shift p q) ((t.sub_shift p q).symm cast r)
      = t.shift p ((t.sub p).2.shift q r)
\end{lstlisting}

\paragraph{Haskell.}
\begin{lstlisting}[style=haskell]
-- D3: shift t p Here = p
-- D4: sub (shift t p q) = sub (sub t p) q
-- D5: shift (shift t p q) r = shift t p (shift (sub t p) q r)
\end{lstlisting}


%===================================================================
\section{The BOM directed container}
%===================================================================

\subsection{Definition: \lean{bom}}

\begin{defn}[BOM directed container]
The directed container \lean{bom}$(P, G)$ has:
\begin{align*}
  \mathrm{Shape} &= \textstyle\sum_{v : G.V} \lean{Bom}(P, G, v), \\
  \mathrm{Pos}(v, t) &= \lean{Occ}(t), \\
  \mathrm{root}(v, t) &= \lean{here}, \\
  \mathrm{sub}((v, t), p) &= t.\lean{sub}(p), \\
  \mathrm{shift}((v, t), p, q) &= t.\lean{shift}(p, q).
\end{align*}
This is exactly Construction~\texttt{con:bom}'s directed container of rooted
trees, with labels, sites, and routes riding along the occurrences.
\end{defn}

\paragraph{Lean.}
\begin{lstlisting}[style=lean]
def bom (P : Type) (G : Digraph) : DirectedContainer where
  Shape := (v : G.V) * Bom P G v
  Pos s := Bom.Occ s.2
  root _ := .here
  sub s p := s.2.sub p
  shift s p q := s.2.shift p q
  d1 _ := rfl
  d4 s p q := s.2.sub_shift p q
  d2 _ _ := rfl
  d3 s p := s.2.shift_root p
  d5 s p q r := s.2.shift_assoc p q r
\end{lstlisting}

\paragraph{Haskell.}
\begin{lstlisting}[style=haskell]
-- The BOM DC: shapes are rooted BOM trees (with facility),
-- positions are occurrences.
data BomDC p v e = BomDC
  { bomShape :: (v, Bom p v e)
  , bomPos   :: Occ
  }
\end{lstlisting}


%===================================================================
\section{The sourcing functor}
%===================================================================

% --- Bom.routeOf --------------------------------------------------
\subsection{Definition: \lean{Bom.routeOf}}

\begin{defn}[Inbound route of an occurrence]
The \textbf{inbound route} $\lean{routeOf}(t, p)$ of occurrence $p$ in BOM $t$
at facility $v$ is a path in $G^{\mathrm{op}}$ from $v$ to the sub-BOM's
facility, built by appending the routes along the path from root to $p$:
\begin{itemize}
  \item $\lean{routeOf}(t, \lean{here}) = \mathrm{nil}_{G^{\mathrm{op}}}(v)$ (empty inbound route);
  \item $\lean{routeOf}(\lean{node}(\ldots, \mathrm{routes}, \mathrm{children}),
    \lean{inside}(i, p)) = \mathrm{routes}(i) \mathbin{+\!\!+} \lean{routeOf}(\mathrm{children}(i), p)$.
\end{itemize}
\end{defn}

\paragraph{Lean.}
\begin{lstlisting}[style=lean]
def routeOf : {v : G.V} -> (t : Bom P G v) -> (p : Occ t) ->
    Digraph.Path G.op v (t.sub p).1
  | v, _, .here => Digraph.Path.nil (G := G.op) v
  | _, .node _ _ _ _ routes children, .inside i p =>
      (routes i).append ((children i).routeOf p)
\end{lstlisting}

\paragraph{Haskell.}
\begin{lstlisting}[style=haskell]
routeOf :: Bom p v e -> Occ -> [e]
routeOf _ Here = []
routeOf (Node _ _ children) (Inside i p) =
  route (children !! i) ++ routeOf (thd (children !! i)) p
  where route (_, r, _) = r
\end{lstlisting}


% --- routeOf_shift -----------------------------------------------
\subsection{Theorem: \lean{Bom.routeOf\_shift} (Consolidation)}

\begin{thm}[Consolidation --- F2 of the sourcing functor]
For all $t : \lean{Bom}(P, G, v)$, $p : \lean{Occ}(t)$, and $q :
\lean{Occ}(\lean{sub}(t, p).\mathbf{2})$:
\begin{align*}
  &\bigl(\lean{sub}(t, \lean{shift}(t, p, q)).\mathbf{1},\;
    \lean{routeOf}(t, \lean{shift}(t, p, q))\bigr) \\
  &\quad=\;
  \bigl((\lean{sub}(t, p).\mathbf{2}.\lean{sub}(q)).\mathbf{1},\;
    \lean{routeOf}(t, p) \mathbin{+\!\!+}
    \lean{sub}(t, p).\mathbf{2}.\lean{routeOf}(q)\bigr).
\end{align*}
The route of a concatenated occurrence is the concatenation of the routes.
Operationally: ore reaches the rolling mill \emph{only via} the reduction
plant and the melt shop.  This is the key coherence law (F2) for the sourcing
functor.
\end{thm}

\paragraph{Lean.}
\begin{lstlisting}[style=lean]
theorem routeOf_shift : {v : G.V} -> (t : Bom P G v) -> (p : Occ t) ->
    (q : Occ (t.sub p).2) ->
    (<|(t.sub (t.shift p q)).1, t.routeOf (t.shift p q)|>
        : (u : G.V) * Digraph.Path G.op v u)
      = <|((t.sub p).2.sub q).1,
          (t.routeOf p).append ((t.sub p).2.routeOf q)|>
\end{lstlisting}

\paragraph{Haskell.}
\begin{lstlisting}[style=haskell]
-- consolidation:
-- routeOf t (shift t p q) = routeOf t p ++ routeOf (sub t p) q
propConsolidation :: Bom p v e -> Occ -> Occ -> Bool
propConsolidation t p q =
  routeOf t (shift t p q) == routeOf t p ++ routeOf (subBomTree t p) q
  where subBomTree t p = snd (subBom t p)
\end{lstlisting}


% --- sourcing -------------------------------------------------------
\subsection{Definition: \lean{sourcing}}

\begin{defn}[Sourcing functor (Construction \texttt{con:sourcing})]
The \textbf{sourcing functor} is a covariant directed-container functor
\[
  \lean{sourcing}(P, G) :
  \lean{DContFunctor}(\lean{bom}(P, G),\; \mathcal{N}(G^{\mathrm{op}}))
\]
with:
\begin{itemize}
  \item on shapes: $(v, t) \mapsto v$ (the producing facility of the root product);
  \item on positions: $(v, t, p) \mapsto (\lean{sub}(t, p).\mathbf{1},\;
    \lean{routeOf}(t, p))$ (the component's facility and its inbound route).
\end{itemize}
Laws F0 ($\mathrm{sub}$ compatibility) and F1 (root preservation) hold by
\lean{rfl}; F2 (shift compatibility) is \lean{routeOf\_shift}.
\end{defn}

\paragraph{Lean.}
\begin{lstlisting}[style=lean]
def sourcing (P : Type) (G : Digraph) :
    DContFunctor (bom P G) (Digraph.network G.op) where
  onShapes s := s.1
  onPos s p := <|(s.2.sub p).1, s.2.routeOf p|>
  f0 _ _ := rfl
  f1 _ := rfl
  f2 s p q := s.2.routeOf_shift p q
\end{lstlisting}

\paragraph{Haskell.}
\begin{lstlisting}[style=haskell]
sourcing :: DContFunctor (BomDC p v e) (Network v e)
sourcing = DContFunctor
  { onShapes = \(v, _) -> v
  , onPos    = \(_, t) p -> (fst (subBom t p), routeOf t p)
  -- F0, F1 by definition; F2 is consolidation
  }
\end{lstlisting}


%===================================================================
\section{UAE steel instance}
%===================================================================

\subsection{Products, sites, and graph}

\begin{defn}[UAE steel supply chain]
The UAE steel supply chain (\S6 of the paper) has:
\begin{itemize}
  \item products: \lean{orePellets}, \lean{dri}, \lean{scrap}, \lean{billet},
    \lean{rebar};
  \item sites: \lean{oreBerth}, \lean{reductionPlant}, \lean{meltShop},
    \lean{rollingMill};
  \item legs: \lean{ob2rp} (ore berth $\to$ reduction plant), \lean{rp2ms}
    (reduction plant $\to$ melt shop), \lean{ms2rm} (melt shop $\to$ rolling mill).
\end{itemize}
The shipping graph is a directed path of length~3.
\end{defn}

\paragraph{Lean.}
\begin{lstlisting}[style=lean]
inductive Product
  | orePellets | dri | scrap | billet | rebar

inductive Site
  | oreBerth | reductionPlant | meltShop | rollingMill

inductive Leg
  | ob2rp | rp2ms | ms2rm

@[reducible] def graph : Digraph where
  V := Site; E := Leg
  src := Leg.src; tgt := Leg.tgt
\end{lstlisting}

\paragraph{Haskell.}
\begin{lstlisting}[style=haskell]
data Product = OrePellets | Dri | Scrap | Billet | Rebar
data Site    = OreBerth | ReductionPlant | MeltShop | RollingMill
data Leg     = Ob2rp | Rp2ms | Ms2rm
\end{lstlisting}


\subsection{Sanity checks (examples checked by \lean{rfl})}

\begin{prop}[Sourcing maps steel shape to its facility]
$({\lean{sourcing}}(P, G)).\lean{onShapes}(\lean{steelShape}) = \lean{rollingMill}$
--- the rebar BOM is produced at the rolling mill.
\end{prop}

\begin{prop}[Ore occurrence inbound route]
The ore-pellet occurrence in the rebar BOM is sourced by the three-leg inbound
route: rolling mill $\leftarrow$ melt shop $\leftarrow$ reduction plant $\leftarrow$
ore berth.
\end{prop}

\begin{prop}[Scrap inbound route]
The scrap occurrence is sourced by the one-leg inbound route: rolling mill
$\leftarrow$ melt shop (scrap is received at the melt shop, not shipped in from
a remote site).
\end{prop}

\begin{prop}[Consolidation instance]
Sourcing the concatenated occurrence billet-then-ore equals shifting the
billet's inbound route by the ore's inbound route in the opposite network:
\begin{align*}
  &\lean{sourcing}.f\bigl(\lean{steelShape},\;
    \lean{shift}(\lean{rebarBom}, \lean{billetOcc}, \lean{oreInBillet})\bigr) \\
  &\quad=\;
  \mathcal{N}(G^{\mathrm{op}}).\lean{shift}\bigl(\lean{rollingMill},\;
    \lean{sourcing}.f(\lean{steelShape}, \lean{billetOcc}),\;
    \lean{sourcing}.f(\lean{billetBom}, \lean{oreInBillet})\bigr).
\end{align*}
All checked by \lean{rfl}.
\end{prop}

\paragraph{Lean.}
\begin{lstlisting}[style=lean]
example : (sourcing Product graph).onShapes steelShape = rollingMill := rfl

example : (sourcing Product graph).onPos steelShape oreOcc
    = <|oreBerth, (inb ms2rm).append ((inb rp2ms).append
          ((inb ob2rp).append
           (Digraph.Path.nil (G := graph.op) oreBerth)))|> := rfl

example : (sourcing Product graph).onPos steelShape scrapOcc
    = <|meltShop, (inb ms2rm).append
          (Digraph.Path.nil (G := graph.op) meltShop)|> := rfl

example : (sourcing Product graph).onPos steelShape
      (rebarBom.shift billetOcc oreInBillet)
    = (Digraph.network graph.op).shift rollingMill
        ((sourcing Product graph).onPos steelShape billetOcc)
        ((sourcing Product graph).onPos <|meltShop, billetBom|>
           oreInBillet) := rfl
\end{lstlisting}

\paragraph{Haskell.}
\begin{lstlisting}[style=haskell]
-- sourcingShape rebarShape = RollingMill
-- sourcingPos rebarShape oreOcc
--   = (OreBerth, [Ms2rm, Rp2ms, Ob2rp])  (reversed legs: inbound)
-- sourcingPos rebarShape scrapOcc
--   = (MeltShop, [Ms2rm])
-- consolidation instance: rfl in Lean
\end{lstlisting}


%===================================================================
\section{Translation conventions}
%===================================================================

\begin{enumerate}
\item \textbf{Indexed inductive type.}  \lean{Bom P G v} is indexed by the
  producing facility $v$.  In Haskell, $v$ is a type parameter; the indexing
  guarantee is lost.

\item \textbf{Opposite paths as inbound routes.}  A $G^{\mathrm{op}}$-path from
  $v$ to $u$ is physically a $G$-path from $u$ to $v$ (same edge list, swapped
  endpoints).  In Haskell, we represent both as the same list type and note the
  direction convention.

\item \textbf{Routes on shapes (formalization caveat).}  The paper's BOM shapes
  are bare product-labelled trees; here, they carry routing choices.  See the
  module docstring for why this is necessary.

\item \textbf{Proof bodies omitted.}  Only definitions and theorem
  \emph{statements} are translated.  Lean checks the proofs; the human
  verifier checks that the \emph{statements} match the intended mathematics.
\end{enumerate}

\end{document}
